% Generated draft: requires mathematical and source review.
\documentclass[11pt]{article}
\usepackage[margin=25mm]{geometry}
\usepackage{fontspec}
\setmainfont{Latin Modern Roman}
\usepackage{amsmath,amssymb,mathtools,unicode-math}
\setmathfont{Latin Modern Math}
\usepackage{enumitem,needspace,xurl,hyperref}
\setlength{\parindent}{0pt}
\setlength{\parskip}{5pt}
\setlength{\emergencystretch}{3em}
\newcommand{\N}{\mathbb{N}}
\newcommand{\Z}{\mathbb{Z}}
\newcommand{\Q}{\mathbb{Q}}
\newcommand{\R}{\mathbb{R}}
\newcommand{\C}{\mathbb{C}}
\newcommand{\F}{\mathbb{F}}
\newcommand{\A}{\mathbb{A}}
\newcommand{\PP}{\mathbb{P}}
\newcommand{\E}{\mathbb{E}}
\newcommand{\eps}{\varepsilon}
\newcommand{\dd}{\,\mathrm{d}}
\begin{document}
\section*{UnsolvedMath ID: 30004433}
\textbf{Source problem code:} catalogue entry
\subsection{Definitions and mathematical statement}
The supplied record concerns one-dimensional long-range percolation at decay exponent \(s=2\). In the usual lattice interpretation, let \(G_\beta\) be a random undirected graph on \(\Z\), with edge probabilities depending on displacement and satisfying
\[
 p_\beta(k):=\mathbb P_\beta\bigl(\{0,k\}\in E(G_\beta)\bigr),\qquad
 \lim_{|k|\to\infty}\frac{p_\beta(k)}{\beta |k|^{-2}}=1,
 \qquad \beta>0.
\]
The record does not specify the complete edge law, independence of distinct edges, probabilities at finite displacements, or the parameter family used to define criticality. Thus this notation describes the conventional interpretation of the stated tail, rather than a fully specified probability model.

For a connected, locally finite, infinite graph \(C\), a ray is a one-way infinite simple path. Two rays determine the same \emph{end} if, after deletion of any finite vertex set, tails of both rays lie in the same connected component. Write \(e(C)\) for the number of ends. Equivalently, for the distinctions requested here,
\[
 e(C)=2\quad\Longleftrightarrow\quad
 \sup_{\substack{F\subseteq V(C)\\F\text{ finite}}}N_\infty(C\setminus F)=2,
\]
where \(N_\infty\) counts infinite connected components; \(C\) has infinitely many ends precisely when this supremum is unbounded. No particular infinite cardinality of ends is asserted.

If \(\beta\) is the intended parameter of a specified monotone family, write
\[
 \beta_c:=\inf\left\{\beta>0:
 \mathbb P_\beta\bigl(G_\beta\text{ has an infinite connected component}\bigr)>0\right\}.
\]
The source asks whether both of the following assertions hold at \(s=2\):
\begin{enumerate}
\item At criticality, the infinite cluster has exactly two ends.
\item In the supercritical regime, the infinite cluster has infinitely many ends.
\end{enumerate}
For a specified model having an almost surely unique infinite cluster \(C_\infty\) at the relevant parameters, the intended almost-sure assertions would read
\[
 \mathbb P_{\beta_c}\bigl(e(C_\infty)=2\bigr)=1,
 \qquad
 \mathbb P_\beta\bigl(C_\infty\text{ has infinitely many ends}\bigr)=1
 \quad\text{for every supercritical }\beta.
\]
These displays require existence and uniqueness of \(C_\infty\), including at criticality. The supplied record does not establish those hypotheses, specify whether the critical assertion is conditional on existence, or explicitly give the almost-sure quantifier. If more than one infinite cluster can occur, it also does not specify which cluster, or whether every infinite cluster, is intended. These qualifications remain unresolved; the tail asymptotic alone does not supply the missing model specification.
\subsection{Short English statement}
For one-dimensional long-range percolation with edge probabilities asymptotic to \(\beta |k|^{-2}\), the supplied question asks whether the critical infinite cluster has exactly two ends and a supercritical infinite cluster has infinitely many ends. The full model, the meaning of criticality, and the existence, uniqueness, and probabilistic quantifiers for the infinite cluster require clarification.
\subsection{Sources}
\begin{itemize}
\item Supplied local UnsolvedMath registry record: permanent ID \(30004433\), source problem code \texttt{catalogue entry}, title \emph{Ends of Critical Long-Range Percolation Clusters}, collection \emph{Oberwolfach Reports Open Problems}. No individual proposer is supplied; the proposed year is recorded as 2020.
\item Supplied report citation: \emph{Mini-Workshop: One-sided and Two-sided Stochastic Descriptions} (2021), \emph{Oberwolfach Reports}, pp.\ 601--637, DOI \texttt{10.4171/owr/2020/11}. Supplied source link: \url{https://doi.org/10.4171/owr/2020/11}. The background describes these reports as documenting mathematical workshops and their open-problem discussions.
\item Supplied literature reference, labeled \texttt{OWR-2020-11}: \url{https://ems.press/content/serial-article-files/46847}.
\item Literature assessment dated \texttt{2026-08-22}: the supplied curation labels the problem open, reports that both assertions are posed as open problems in the primary report, and states that no general resolution was identified. This is supplied dated curation, not an independently verified literature assessment; time-sensitive status and cited claims require independent verification.
\end{itemize}
\end{document}
